\abstract % Структурный элемент: РЕФЕРАТ

\keywords{МУЛЬТИ-РОБОТИЗИРОВАННЫЕ СИСТЕМЫ, MAPF, MRTA, ПОИСК БЕСКОНФЛИКТНЫХ МАРШРУТОВ, РАСПРЕДЕЛЕНИЕ ЗАДАЧ, СРАВНИТЕЛЬНЫЙ АНАЛИЗ АЛГОРИТМОВ, ПРОГРАММНЫЙ КОМПЛЕКС}

Объектом разработки является кроссплатформенный программный комплекс для оценки эффективности алгоритмов управления мульти-роботизированными системами.

Цель работы~--- создание программного комплекса, позволяющего проводить тестирование и сравнительный анализ алгоритмов поиска бесконфликтных маршрутов (MAPF) и распределения задач (MRTA) при различных конфигурациях среды.

Методы проведения работы: анализ и классификация алгоритмов на основе научной литературы; объектно-ориентированное проектирование с применением паттерна <<Стратегия>> для обеспечения модульности и расширяемости; реализация алгоритмов на Python с графическим интерфейсом на Pygame; автоматизированное пакетное тестирование с варьированием параметров и статистической обработкой результатов.

В ходе работы реализованы восемь MAPF-алгоритмов (CBS, CBS-SIPP, ECBS, EECBS, M*, CA*, WHCA*, PIBT) и пять MRTA-алгоритмов (венгерский, жадный, последовательный и комбинаторный аукционы, метод минимальной стоимости потока). Разработан графический интерфейс с редактором карт, визуализацией симуляции и режимом пакетного тестирования с автогенерацией графиков. Обеспечена кроссплатформенная сборка в автономный исполняемый файл для Windows, Linux и macOS.

Комплекс устраняет разрозненность существующих инструментов, объединяя алгоритмы MAPF и MRTA в единой среде с визуализацией и пакетным тестированием. Это обеспечивает обоснованный выбор алгоритмов по объективным метрикам. Комплекс применим в научных исследованиях по координации групп агентов, прототипировании систем складской логистики, автономного транспорта и спасательных операций.